\chapter{Předpisy a nařízení}
V případě Open Data je potřeba se seznámit s platnou legislativou upravující tuto oblast.

Zákonem upravující otevřená data je zákon 106/1999 Sb. o svobondém přístupu k informacím.  Tento zákon určuje povinnosti především pro subjekty, které jsou povinné se tímto zákonem řídit. Dobrovolní zveřejňovatelé se ale musejí tímto zákonem řídít také.

Otevřená data by podle portálu opendata.gov.cz měla splňovat tyto podmínky \parencite{CoJsouOtevrena2022}:
\begin{enumerate}[]
	\item{Přístupná jako datové soubory ve strojově čitelném a otevřeném formátu s úplným a aktuálním obsahem databáze nebo agregovanou statistikou} 
	\item{Opatřená neomezujícími podmínkami užití}
	\item{Evidovaná v Národním katalogu otevřených dat (NKOD) jako přímé odkazy na datové soubory}
	\item{Opatřená dokumentací}
	\item{Dostupná ke stažení bez technických překážek (registrace, omezení počtu přístupů, CAPTCHA, apod.)}
	\item{Připravena s cílem co nejsnazšího strojového zpracování programátory apod.}
	\item{Opatřená kontaktem na kurátora pro zpětnou vazbu (chyby, žádost o rozšíření, apod.)}
\end{enumerate}
